\subsection{Телеметрия и работа с датчиками}

У Pioneer есть множество датчиков. Мы научимся получать от них данные.

\subsubsection*{Основные датчики}
\begin{itemize}
    \item \textbf{Барометр:} Измеряет высоту.
    \item \textbf{IMU (Гироскоп + Акселерометр):} Измеряет углы наклона и ускорения.
    \item \textbf{Дальномер (ToF):} Лазерный датчик расстояния (смотрит вниз), измеряет дистанцию до пола.
    \item \textbf{Система оптического потока (OptiFlow):} Камера, смотрящая вниз, для удержания позиции.
\end{itemize}

\subsubsection*{Чтение данных дальномера}
Дальномер полезен для точного удержания высоты при посадке или полете над препятствиями.

\begin{codefile}[python]{python/examples/get_telemetry/GET_dist_sensor_data.py}
\end{codefile}

\subsubsection*{Локальная позиция}
Дрон постоянно вычисляет свои координаты $(x, y, z)$ относительно точки старта. Эти данные можно получить через \texttt{get\_local\_position\_lps()}.

\begin{taskbox}[Черный ящик]
Напишите программу, которая во время полета каждые 0.5 секунды записывает в текстовый файл \texttt{flight\_log.txt} текущее время, высоту и заряд батареи.
\end{taskbox}

\begin{taskbox}[Мягкая посадка]
Используя дальномер, реализуйте алгоритм мягкой посадки: дрон снижается со скоростью 0.5 м/с. Когда высота становится меньше 0.3 метра, скорость снижения уменьшается до 0.1 м/с.
\end{taskbox}
